\documentclass{article} % For LaTeX2e
\usepackage{iclr2027_conference,times}

\input{math_commands.tex}

\usepackage{hyperref}
\usepackage{url}
\usepackage{booktabs}
\usepackage{graphicx}
\usepackage{amsmath}

\title{Capability Atoms: Shaping What a Distilled\\Code Agent Can Do from a Few-Shot Task Specification}

% Anonymized for submission.
\author{Anonymous authors\\
\And
\\
}

\newcommand{\fix}{\marginpar{FIX}}
\newcommand{\new}{\marginpar{NEW}}
% Working-draft marker; remove for submission.
\newcommand{\todo}[1]{\textcolor{red}{[TODO: #1]}}
\usepackage{xcolor}

\iclrfinalcopy % TODO(submission): remove for anonymized submission
\begin{document}

\maketitle

\begin{abstract}
%% v0.1 (2026-08-10). Numbers marked "pilot" are single-seed; the (size,
%% budget) sweep and agent-environment results will replace them as they land.
Deploying a large language-model agent often means distilling it into a small
model that must handle one task domain, specified only by a handful of example
queries. Deciding which requests fall inside that domain, and training a
student whose competence ends where the domain ends, are usually treated as
separate problems solved with semantic similarity---which cannot work when two
tasks share their surface text but differ in the skills they require. We
represent both problems in one space: the gradient a trajectory induces in the
student encodes what learning that task would change, and a sparse dictionary
over these gradients yields \emph{capability atoms}---interpretable units that
support task boundaries, capability measurement, and data selection in a
single representation. On a hard pair sharing identical questions but
differing in action space (Python vs.\ chain-of-thought), gradient boundaries
reach AUROC $1.000$ where trajectory embeddings reach $0.687$ (pilot); atom
support resolves within-domain skills (SQL-join, AUROC $0.906{\pm}0.025$)
where dense subspaces fail ($0.615$); boundaries carry finite-sample conformal
coverage guarantees and are stable across nine reference models from 0.8B to
27B parameters. Boundary-guided selection matches the best in-domain execution
accuracy while sharpening the student most strongly away from out-of-domain
capability, and a refusal module separates in- from out-of-domain requests
with 100\%/0\% refusal rates at a 6.7-point cost (pilot). We report our
failed pre-registered gate for a gradient-norm capability proxy, and the
mechanism analysis that replaced it.
\end{abstract}

\section{Introduction}
%% Skill check: para-1 concrete failure; para-2 crux; para-3 insight; para-4 claims.

A team builds an agent on a frontier model that files expense reports by
writing Python against internal APIs. The model is too expensive to serve, so
they distill a 2B student from the teacher's trajectories. Two questions
decide whether the student is deployable, and neither is answered by current
distillation practice: given a new request, is it \emph{inside} the job the
student was built for; and does the student's competence in fact end where
that job ends, rather than silently covering less---or more---than intended?

The obvious tool for both questions is semantic similarity over the request
text. It is structurally the wrong tool. Consider the same grade-school math
question solved twice: once as a Python function, once as chain-of-thought
prose. The requests are \emph{identical strings}, so no function of the query
text can separate them; yet they define different tasks, exercising different
skills, and a student distilled for one does not acquire the other
(Sec.~\ref{sec:boundary}). Task identity lives in what a model would have to
\emph{learn}, not in what the request \emph{says}.

Our starting point is to represent tasks by their learning demands: the
gradient $g(q)$ that a teacher trajectory for query $q$ induces at the
student's initialization is a direct encoding of ``which parameters must move
to learn this''. A sparse dictionary over trajectory-block gradients then
factors these demands into \emph{capability atoms}: recurring directions that
align with recognizable skills (the atoms most discriminative for SQL-join
tasks activate precisely on \texttt{JOIN \ldots ON} clauses). One
representation supports the full pipeline: a task specification is the atom
support of $k$ example queries with a conformal acceptance radius; a
student's capability is the supply each atom received from its training data;
data selection maximizes supply to the specified atoms under a token budget.

\textbf{Contributions.} (i)~We formalize \emph{distill-to-specification}:
producing, from a teacher and $k$ example queries, both a boundary predictor
with finite-sample coverage guarantees and a student whose measured
capability matches the specified domain on both sides
(Sec.~\ref{sec:setting}). (ii)~We show task boundaries must be estimated in
learning space, not semantic space: on identical-query hard pairs, gradient
features reach AUROC $1.000$ vs.\ $0.687$ for trajectory embeddings, saturate
at $k{=}5$ examples, and are stable across nine reference models spanning
0.8B--27B, base and instruct (Sec.~\ref{sec:boundary}, pilot).
(iii)~We introduce capability atoms and show a single sparse representation
covers both topic-level boundaries (AUROC $0.94$--$1.00$) and within-domain
skill resolution ($0.906$ vs.\ $0.615$ for dense subspaces)
(Sec.~\ref{sec:atoms}). (iv)~We report a disciplined failure: our
pre-registered gate for a gradient-norm capability proxy failed twice
($\rho=-0.40, -0.30$ vs.\ threshold $-0.5$); the diagnosis (style drift and
exposure bias) motivates measuring capability as predicted atom supply,
validated against execution outcomes (Sec.~\ref{sec:capability}).
(v)~In budget-matched comparisons, boundary-guided selection matches the best
in-domain execution accuracy (90\%, pilot) while most strongly removing
out-of-domain capability, and we show mechanistically that this sharpening is
supply starvation rather than gradient conflict: fewer than 15\% of
cross-domain gradient inner products are negative, while kernel mass predicts
where loss falls ($\rho=0.53$, $p=0.04$) (Sec.~\ref{sec:distill}).

\section{Related work}\label{sec:related}
\todo{Compress novelty-scan clusters into 5 paragraphs: agent distillation
(SmartAD, code-tool distillation, on-policy lines) optimizes benchmark scores
without task specifications; gradient-based selection (LESS line) and
attribution (TRAK, Gradient Atoms) diagnose rather than design; conformal OOD
rejection operates at serving time, decoupled from training; learnware
retrieves models matching capability specifications where we synthesize them;
unlearning removes capability post hoc where we shape it during distillation.}

\section{Problem setting}\label{sec:setting}
A teacher agent $\pi_T$ acts by emitting executable code; a rollout on query
$q$ yields a trajectory $\tau(q)$ whose success is decided by an execution
verifier. A task domain is a distribution $P_\mathcal{T}$ over queries given
only through $k$ examples $Q_k \sim P_\mathcal{T}$, with $k$ between 5 and a
few hundred. From $(\pi_T, Q_k)$ and a base student $\theta_0$ we must
produce: (1) a boundary predictor $\hat{B}(q) \in \{\text{in},\text{out}\}$
with distribution-free coverage $\Pr_{q\sim P_\mathcal{T}}[\hat{B}(q)=
\text{in}] \ge 1-\alpha$; and (2) a student $\theta_S$ evaluated on the
two-sided objective: coverage risk $R_{\mathrm{cov}}$ (in-domain requests the
student fails) and leakage $L$ (out-of-domain requests the student performs).
We report $(R_{\mathrm{cov}}, L)$ frontiers under matched token budgets and,
across student sizes, the minimal capacity achieving target coverage.

Capability and behavior are measured separately throughout: distillation can
flip an agent's action space within tens of optimizer steps while transferring
only part of the underlying competence (Sec.~\ref{sec:distill}), so any
metric reading only output format overestimates capability transfer.

\section{Method}\label{sec:method}
%% Scope-level per PI instruction; each part carries its evidence pointer and
%% will be updated as results land. Notation kept minimal.

\paragraph{Learning-demand features.}
For query $q$ with trajectory $\tau(q)$, the feature is
\begin{equation}
g(q) \;=\; \mathrm{normalize}\!\left(P\, \nabla_\theta\,
\ell(\tau(q);\theta_0)\right),
\end{equation}
the (sketched) gradient of the trajectory's action-token loss at the student
base. With zero-initialized LoRA adapters, the adapter-$B$ gradient equals
$G A^{\!\top}$---a free random projection of the full gradient $G$---and a
per-module Johnson--Lindenstrauss step reduces each sample to $\sim$12k
dimensions. Geometry is preserved (a 768-d control loses nothing measurable),
and the resulting structure is stable across nine reference models
(0.8B--27B, base and instruct, two generations).

\paragraph{Capability atoms.}
Trajectories decompose into blocks (statements); block gradients decompose
sparsely over a learned dictionary $\{a_1,\dots,a_m\}$
(mini-batch dictionary learning; $m{=}128$, $\sim$5 active atoms per block).
Atoms are operational skill units: the atoms separating SQL-join from
plain-SQL tasks activate on \texttt{JOIN} clauses. The identity
$g(q) = \sum_t g_t(q)$ makes granularity an aggregation choice within one
representation: whole-trajectory sums resolve topics; block-level support
resolves skills. \todo{stability analysis of the dictionary across seeds;
AST-level blocks for intra-line skills (pandas case).}

\paragraph{Boundary with conformal coverage.}
The specification of $Q_k$ is its atom support (equivalently, at topic scale,
the spectral subspace of $\{g(q_i)\}$). A membership score is calibrated by
split conformal prediction: with $n_{\mathrm{cal}}$ held-out examples, the
acceptance threshold guarantees $1-\alpha$ coverage of in-domain queries.
Measured coverage tracks the nominal level (93--97\% at $\alpha=0.1$, pilot).
The boundary is deliberately inclusive---its job is to cover the intended
topic, and precision at skill level is delegated to atom support.

\paragraph{Capability as supplied mass, verified behaviorally.}
We do not measure student capability by post-hoc gradient norms: that proxy
failed its pre-registered gate (pooled Spearman $-0.40$ and $-0.30$ against a
$-0.5$ threshold, two metric variants), with failures localized to style
drift (residuals against teacher references rise as the student forms its own
style) and exposure bias (teacher-forced scores cannot see decoding failure).
Instead, capability is \emph{predicted forward}: training on data $D$
supplies each atom with mass $\sum_{x\in D}\langle g(q), g(x)\rangle$, and
this prediction is validated against execution outcomes. Kernel mass ranks
measured loss changes across domains ($\rho=0.53$, $p=0.04$, pilot);
per-student validation on budget-matched cohorts is pre-registered at
AUROC $\ge 0.75$. \todo{results of the cohort validation (E2-v3).}

\paragraph{Boundary-guided distillation and shaping control.}
Selection ranks a heterogeneous trajectory pool by specification score and
fills a token budget (a first-order special case of maximizing atom coverage,
which is submodular). Out-of-domain leakage is controlled by \emph{supply
allocation} rather than opposition: fewer than 15\% of cross-domain gradient
inner products are negative, so narrowing the diet starves unspecified
capability---predictably---while parameters move for the specified task.
An optional refusal module converts the lowest-scoring pool items into
refusal targets, shaping behavior on the same boundary
(100\% out-of-domain refusal at 0\% in-domain over-refusal, pilot).

\section{Experiments}\label{sec:experiments}
%% Each subsection opens with its falsifiable question (skill rule).

\subsection{Can boundaries be estimated where semantics fail?}\label{sec:boundary}
\todo{E1 full: hard-pair table (grad 1.000 / emb-traj 0.687 / emb-query 0.404
$\approx$ chance by construction), $k$-efficiency (saturation at $k{=}5$;
embeddings degrade with $k$), dimension- and template-controls, nine-model
robustness matrix, PCA figure.}

\subsection{Do atoms resolve skills within a domain?}\label{sec:atoms}
\todo{Dense-vs-sparse table: SQL-join 0.906$\pm$0.025 vs.\ 0.615 (dense) and
0.866 (embedding, surface-advantaged by construction); coarse-scale parity
(0.94--1.00); atom interpretability panel; dictionary stability; pandas
limitation (intra-line skills need finer blocks).}

\subsection{Does measured capability track behavior?}\label{sec:capability}
\todo{Gate narrative: v1/v2 failures with per-checkpoint analysis (late
within-checkpoint $\rho=-0.70$; lowest-decile success 96\%), style-drift
evidence (per-row residual 15.1 vs.\ 1.3), atom-supply validation on E3
cohorts (pre-registered).}

\subsection{Does boundary-guided distillation shape capability?}\label{sec:distill}
\todo{Budget-matched conditions (all/random/embedding/gradient/refusal):
in-domain execution, action-space capture (63\%$\to$0\% CoT in every
condition), out-of-domain loss movement, refusal split; (size, budget)
coverage surface from the 39-run sweep; supply-starvation mechanism (M1).}

\subsection{Agent environments}\label{sec:agents}
\todo{AppWorld / BFCL-v4 / $\tau^2$-bench topic-scoped results (tier 2);
teacher pool from three current open teachers with execution-verified
quality gates.}

\section{Limitations and pre-registration}\label{sec:limitations}
Pilot-scale numbers are single-seed and use verified pseudo-trajectories or
two teachers; sweep-scale results carry seeds and confidence intervals as
they land. The capability-proxy gate governs what may enter the training
loop: geometric readings stay outside it until the pre-registered validation
passes. Suppression claims are behavioral unless paired with unlearning;
base-model prior capability cannot be removed by distillation alone, and our
base/instruct dual track separates what shaping adds from what priors
contribute. Dictionary hyperparameters materially affect atom quality (an
over-sparse setting fails); stability analyses are required before
skill-level claims generalize. \todo{keep synchronized with amendment log.}

\subsubsection*{Reproducibility statement}
All experiments run from configuration files with sha256-stable splits; every
result row in this paper maps to a line in the experiment log and a commit.
Code, logs, and the amendment history will be released.

\bibliography{iclr2027_conference}
\bibliographystyle{iclr2027_conference}

\appendix
\section{Appendix}
\todo{full protocols; verifier sandbox details; teacher prompts and
anti-degeneracy filters; per-run manifests; additional figures.}

\end{document}
